\documentclass[11pt]{article}
\usepackage{fontspec}
\setmainfont{Times New Roman}
\setsansfont{Arial}
\setmonofont{Consolas}
\usepackage{amsmath,amssymb}
\usepackage{geometry}
\usepackage{microtype}
\geometry{a4paper,margin=3cm}
\setlength{\parindent}{1.25em}
\setlength{\parskip}{0pt}
\emergencystretch=4em
\allowdisplaybreaks
\let\A\relax
\newcommand{\A}{\mathfrak A}
\newcommand{\R}{\mathfrak R}

\begin{document}

\begin{center}
\textbf{\S\ 22. \quad Формы 9. поряда (систем $s_{III}$).}
\end{center}

Свијанье $s$ с формами 5. поряда (систем III) и с продуктом $\A\cdot \nu$ (\S\ 10).

\noindent\textit{Незводиме формы:}
\[
\begin{array}{c@{\,}c|c|c|c}
\text{А)}& \cdot & \cdot & (K_{\nu}su)^2 & \bigstar\\[0.45em]
& \cdot & ((k\nu x)^2,s,u)^2,\ K_{\nu}s_{\nu} & ((k\nu x)^2,s,u)^3 & \\[0.45em]
& (k\nu x)^2s_{\nu},\ ((k\nu x)^3,s,u) & ((k\nu x)^3,s,u)^2 & \cdot & \cdot\\[0.45em]
& (k\nu x)^3s_{\nu} & ((k\nu x)^3,s,u)s_{\nu},\ (K_{\nu}(k\nu x)^3,s,u) & \cdot & \cdot\\[0.45em]
& \bigstar\ (K_{\nu}su)^3 & (K_{\nu}su)^4 & \cdot & (K_{\nu}^2su)^4
\end{array}
\]
\[
\begin{array}{c@{\,}c|c@{\quad}c@{\,}c}
\text{Б)}& (j\nu x)s_{\nu},\ ((j\nu x)^2,s,u) & ((j\nu x)^2,s,u)^2 && \\
& ((j\nu x)^3,s,u) & \cdot && \text{Ц)}\quad s_{\nu}(\A su),\ (\A_{\nu}su),
\end{array}
\]
\[
\begin{array}{c@{\quad}c|c|c}
\text{D)}& \cdot & ((aHu),s,u)^2,\ ((aHu)^2,s,u) & \bigstar\\[0.45em]
& (aHu)s_\eta,\ (a_\eta su) & (aHu)^2s_\eta,\ (a_\eta su)^2 & \\[0.45em]
& \cdot & (a_\eta^2su) & \\
\bigstar& ((aHu),s,u)^3,\ ((aHu)^2,s,u)^2 & ((aHu),s,u)^4 & \\
& (a_\eta su)^3,\ (a_\eta(aHu),s,u)^2,\ (a_\eta(aHu)^2,s,u) & (a_\eta(aHu),s,u)^3 &
\end{array}
\]

\noindent\textit{Редукције:}

А) Редукције следујут просто из редукциј \S\ 21 чрез једнократно свијанье. Ряд форм $K_{\nu}s_{\nu}(Ksu)^2$ има редуценты $a_{\nu}s_{\nu}(asu)^2$ и $\theta_{\nu}s_{\nu}(\theta su)^2$ како факторы. Из того разклад выших форм $K_{\nu}^2(Ksu)^2$, $K_{\nu}^2(Ksu)^3$ по наставку:
\[
0=a_{\nu}s_{\nu}(asu)^2(a\theta u)=K_{\nu}s_{\nu}(Ksu)^2=\theta_{\nu}s_{\nu}(\theta su)^2(a\theta u).
\]

Б) Ряд форм $(j\nu x)^2s_{\nu}$: редуцент $a_{\nu}^2s_{\nu}$ из
\[
(j\nu x)^2s_{\nu}=3a_{\nu}^2s_{\nu}(a\theta u)^2.
\]
$((j\nu x)^3,s,u)^2$: редуцент $a_{\nu}^2s_{\nu}$ из
\[
((j\nu x)^3,s,u)^2=-6a_{\nu}^2s_{\nu}(a\theta u)^2(\theta su).
\]

Ц) Формы $\A_{\nu}(\A su)^2$ и $s_{\nu}(\A su)^2$: редуцент $g_{\nu}$ из формулы (27.)а.

Форма $\A_{\nu}s_{\nu}$: 1) редуцент $a_{\nu}^2s_{\nu}$ из $a_ga_{\nu}^2s_{\nu}=\frac23\A_{\nu}s_{\nu}$;

2) разклада се по формули (27.)а чрез двојну редукцију израза $(\A s\widehat{\nu}x)^2$.

Из того: разклад выше формы $a_\eta s_\eta$.

D) Ряд форм $(aHu)^2s_\eta(asu)$: редуцент $a_{\nu}s_{\nu}(asu)^2$ чрез двојну редукцију ряда форм $[a_{\nu}s_{\nu}(asu)^2]_{\widehat{\nu_1}x}$; редукција $(aHu)^2s_\eta(asu)^2$ из $a_{\nu}s_{\nu}(asu)^2$ и $a_{\nu}s_{\nu}(asu)^3$. Из того редукција $(a_\eta,s,u)^4$.

Ряд форм $a_\eta(aHu)s_\eta$: редуценты $a_{\nu}^2s_{\nu}$, $a_\eta^2s_\eta$ чрез двојну редукцију израза $(\A s\widehat{\nu}x)^3$, бо $a_\eta^2s_\eta$ не наставаје из редуцибилного члена $a_{\nu}^2s_{\nu}(\nu\nu_1x)$ формы $a_\eta(aHu)s_\eta$ чрез свијанье.

Ряд форм $(a_\eta^2su)^2$: редуцент $a_{\nu}^2s_{\nu}$ из
\[
a_{\nu}^2s_{\nu}s_{\nu_1}=-\frac12a_\eta^2(Hsu)^2.
\]
Форма $(a_\eta(aHu),s,u)$ разклада се како трансвекција чрез функционалны детерминант.

Наставајут редукцијске формулы:
\[
\tag{30}
\begin{aligned}
\mathrm{a)}\quad 0&=(aHu)^2(asu)s_\eta-a_\eta(asu)(Hsu)^2-a_\eta(aHu)(asu)(Hsu),\\
\mathrm{b)}\quad 0&=(aHu)^2(asu)^2s_\eta-a_\eta(aHu)(asu)^2(Hsu),\\
\mathrm{c)}\quad 0&=a_\eta(asu)^2(Hsu)^2,\\
0&=a_\eta s_\eta+\frac14a_{\nu}^2\cdot g-\frac14s\cdot a_\eta^2,\\
0&=a_\eta(aHu)s_\eta-\frac12a_\eta^2(Hsu),
\quad 0=a_\eta^2(Hsu)^2,\ldots,
\quad 0=a_\eta^2s_\eta.
\end{aligned}
\]

\noindent\textit{Последки:}

1) $(j\nu x)^2s_{\nu}$, $\A_{\nu}s_{\nu}$, $\A_{\nu}(\A su)^2$ и следовательно $N_{\nu}(Nsu)^2$, $(aHu)^2(asu)s_\eta$, $a_\eta(aHu)s_\eta$, $a_\eta^2(Hsu)^2$ сут редуценты; $a_\eta s_\eta$ јест разкладна форма, ктора при всех једно- и двојкратных свијаньах преходи в разкладне формы.

2) Формы 5. поряда $(5,\lambda)$, $(6,\lambda)$ и т. d. не входят в свијанье с $s^2$.

\end{document}

